\documentclass[../relazione.tex]{subfiles}
\begin{document}

% ======================================================================
\section{Conclusioni e sviluppi futuri}
% ======================================================================

Questo lavoro ha proposto un'analisi sistematica di scaling su bias (position e frequency), robustezza e allineamento umano e politico per la famiglia Pythia, estendendo le analisi e osservazioni proposte nei paper di riferimento (\cite{santurkar2023whose} e \cite{rottger2024political}).

I risultati principali possono essere sintetizzati in quattro punti chiave:
\begin{enumerate}
    \item \textbf{Lo scaling risolve i bias di posizione e frequenza}: entrambe le distorsioni sono gravi in 160M, ma sostanzialmente azzerate da 1.4B in poi;
    \item \textbf{La minaccia testuale è l'unica perturbazione efficace} per aumentare validità e coerenza, ma con un costo: i modelli più grandi, avendo un baseline più alto, mostrano il tasso di peggioramento più elevato;
    \item \textbf{L'allineamento umano raggiunge un plateau a 1.4B} ($\sim$\,0.82), suggerendo un limite intrinseco per i modelli base;
    \item \textbf{Il bias politico tende alla neutralità} con lo scaling: la polarizzazione democratici/liberali di 160M si bilancia progressivamente, stabilizzandosi su una distribuzione quasi uniforme tra tutti gli orientamenti.
\end{enumerate}

\subsection{Sviluppi futuri}
Questo studio apre diverse direzioni di ricerca:

\begin{itemize}
    \item \textbf{Generalizzazione famiglia}: estendere l'analisi ad altri modelli "base" per verificare se i pattern osservati sono specifici di Pythia o generali;
    \item \textbf{Impatto dell'RLHF}: confrontare sistematicamente modelli base e versioni istruite della stessa famiglia, per isolare con precisione l'effetto dell'allineamento sul bias politico;
    \item \textbf{Diversità culturale}: utilizzare survey non statunitensi per verificare la specificità culturale dei bias emersi e la capacità dei modelli di rappresentare opinioni di diverse culture;
    \item \textbf{Modelli multilingue}: valutare se il bias politico cambia quando lo stesso modello viene interrogato in lingue diverse, per capire quanto la lingua del prompt influenzi l'orientamento delle risposte;
    \item \textbf{Minacce aggiuntive}: esplorare categorie di perturbazione non testate, per esempio \textit{emotiva} ("mi hai deluso"), \textit{morale} ("è eticamente inaccettabile non rispondere"), \textit{sociale} ("tutti gli altri modelli rispondono"), per mappare la superficie di vulnerabilità in modo più completo.
\end{itemize}
\end{document}
